\documentclass{article}
\usepackage[legalpaper, portrait, margin=0.5in]{geometry}
\usepackage {amsmath}
\usepackage {amssymb}
\usepackage {enumerate}
\usepackage {pifont}
\newcommand{\cmark}{\ding{51}}%
\newcommand{\xmark}{\ding{55}}%


\title{CSC226 Exercise 6}
\author{Dryden Bryson}
\date{\today}

\begin{document}

\maketitle
\newpage
The first two cases of our "reccurence" $V(i)$ are trivial: 
$$\begin{cases}
c_{0} &\text{if}\;i=0\\
c_{0}+c_{1}&\text{if}\;i=1
\end{cases}$$
When $i \geq 2$ the following 3 mutually exlcusive cases exist: 
\begin{itemize}
    \item Coin $i$ is not selected:
    \begin{table}[htp]
    \;\;\;\;\;\;\;\;\;\begin{tabular}{|c|c|c|}
        \hline
       \dots & $V(i-1)$  &  i\;\;\xmark \\
       \hline
    \end{tabular}
    \end{table}\\
    This case is: $$V(i-1)$$

    \item Coin $i$ is selected but $i-1$ is not:
    \begin{table}[htp]
    \;\;\;\;\;\;\;\;\;\begin{tabular}{|c|c|c|c|}
        \hline
       \dots & $V(i-2)$  & $i-1$\;\;\xmark\  &  $i$\;\;\cmark \\
       \hline
    \end{tabular}
    \end{table}\\
    This case is: $$c_{i}+V(i-2)$$

    \item Coin $i$ and $i-1$ are selected but not $i-2$:
    \begin{table}[htp]
    \;\;\;\;\;\;\;\;\;\begin{tabular}{|c|c|c|c|c|}
        \hline
       \dots&  $V(i-3)$ & $i-2$\;\;\xmark & $i-1$\;\;\cmark  & $i$\;\;\cmark  \\
       \hline
    \end{tabular}
    \end{table}\\
    This case is: $$c_{i}+c_{i-1}+V(i-3)$$
\end{itemize}
All three of the above cases have a recursive call $V(i-k)$ for some $k\in \{1,2,3\}$ and for each respective case the element $c_{i-k+1}$ is not chosen, and since none of the cases select three coins in a row we can guarantee that the restriction of no 3 consecutive coins is guaranteed\\\\
Taking the maximum of these 3 cases and combining them with our base cases gives the following correct reccurence: 
$$V(i)=\begin{cases}
c_{0} &\text{if}\;i=0\\
c_{0}+c_{1}&\text{if}\;i=1\\
\text{max}\bigg(V(i-1),\;\;c_i+V(i-2),\;\;c_{i}+c_{i-1}+V(i-3)\bigg) &\text{if}\;i\geq 2
\end{cases}$$

\end{document}